\section{Conclusioni}
Questo lavoro ha mostrato che una rotonda pedonale può essere interpretata come meccanismo di \textit{Distributed Artificial Intelligence passiva}. L'ordine collettivo non è prodotto da agenti sofisticati o da comunicazione esplicita, ma dalla combinazione di topologia, floor fields e vincoli di occupazione.

L'analisi del codice NetLogo e delle cinque campagne sperimentali consente di trarre quattro conclusioni principali:
\begin{itemize}
\item la geometria della rotonda trasforma crossing flows potenzialmente conflittuali in flussi tangenziali ordinati;
\item l'auto-organizzazione osservata è principalmente geometry-driven, come mostrato dall'invarianza di $\phi$ rispetto a $K_d$;
\item esiste un compromesso critico tra beneficio di separazione e costo geometrico dell'isola centrale, con evidenza di \textit{clogging effect} per raggi troppo grandi;
\item l'eterogeneità psicologica disgrega la compattezza laminare generata dal Floor Field premiando i perfiles aggressivi individualmente ma paralizzando lo scenario globale in micro-shockwaves asincrone.
\end{itemize}

Dal punto di vista metodologico, il lavoro conferma la solidità del Floor Field Model come strumento per studiare crowd dynamics in geometrie complesse. Dal punto di vista concettuale, propone una lettura più forte: l'urban design può essere interpretato come forma di computazione distribuita incorporata.

Gli sviluppi futuri più naturali includono il test di un Threshold di tolleranza (\textit{Critical Mix}) per popolazioni fortemente eterogenee, nonché il confronto con un incrocio privo di isola centrale e una validazione comparativa con social force models continui.
